\documentclass{report}
\usepackage{amsmath,amssymb}
\setlength{\parindent}{0mm}
\setlength{\parskip}{1em}
\begin{document}
\begin{center}
\rule{6in}{1pt} \
{\large Michael Brutz \\
{\bf A Least-Squares Finite Element Method for a Glacier Model Using a Generalized Glen's Flow Law}}

2210 Walnut St \\ Unit 1 \\ Boulder \\ CO 80302
\\
{\tt m.brutz@gmail.com}\\
Chad Westphal\\
Marian Brezina\end{center}

In this talk we present progress towards an efficient and accurate
computational model for non-Newtonain glacier flow. As a part this, we
generalize Glen's Law which is a widely used power law model for the
material. A common criticism of the standard law is it has difficulty
matching data in both low and high stress regimes simultaneously. Our
generalization of the law largely overcomes this difficulty and
furthermore does not suffer from numerical difficulties that have been
observed in the literature for standard Glen's law flow. In regards to
the flow model itself, we focus on a stress-velocity-pressure formulation
that resembles a Stokes system with a nonlinear viscosity. The
discretization we use employs a least squares finite element principle to
minimize the error in an appropriate norm. This produces symmetric
positive definite linear systems that are amenable to efficient multigrid
solvers. Aside from an overview of the model, we also present some
preliminary numerical results for 2-D test problems and discuss current
work.


\end{document}
